\documentclass[11pt]{article}
\usepackage[margin=1in]{geometry}
\usepackage{amsmath,amssymb,amsthm}
\usepackage{graphicx}
\usepackage{mathrsfs}
\usepackage{hyperref}

\newtheorem{theorem}{Theorem}
\newtheorem{lemma}{Lemma}
\newtheorem{remark}{Remark}

\newcommand{\Anu}{\mathcal A_{\nu}}
\newcommand{\Lcal}{\mathcal L}
\newcommand{\R}{\mathbb R}

\title{A Norming Family of Uniform Numerical Bollobas Operators}
\author{Run \texttt{fa\_banach\_001}, agent \texttt{agent\_lane\_04}}
\date{June 23, 2026}

\begin{document}
\maketitle

\section*{Source and Target}

The source paper is S. Dantas, M. Jung, and O. Roldan, \emph{Norm attaining operators which satisfy a Bollobas type theorem}, arXiv:1910.05726. After Proposition 2.8, the authors construct compact operators \(T\in \Anu(H)\), where \(H\) is a separable infinite-dimensional real Hilbert space, such that
\[
1=\nu(T)<\|T\|,
\]
and such that the Bollobas modulus in the definition of \(\Anu(H)\) can be chosen independently of the particular operator \(T\). They write:

\begin{center}
\includegraphics[width=0.95\textwidth]{figures/open_problem_crop.png}
\end{center}

The source does not introduce a separate dual-pairing interpretation for this
sentence.  In the numerical-radius part of the paper, the scalar quantity
attached to an operator is the quadratic form \(x\mapsto \langle Tx,x\rangle\).
We therefore formalize ``norming for the whole space'' in this source-local
sense and prove that the answer is yes.

\section*{Definitions}

For an operator \(T\in\Lcal(H)\), its numerical radius is
\[
\nu(T)=\sup_{\|x\|=1}|\langle Tx,x\rangle|.
\]
In a real Hilbert space, the states in the definition of \(\Anu(H)\) reduce to pairs \((x,x)\), so \(T\in\Anu(H)\) with \(\nu(T)=1\) means: for every \(\varepsilon>0\) there is \(\eta(\varepsilon,T)>0\) such that whenever \(x\in S_H\) and
\[
|\langle Tx,x\rangle|>1-\eta(\varepsilon,T),
\]
there is \(x_1\in S_H\) with \(\langle Tx_1,x_1\rangle=1\) and \(\|x_1-x\|<\varepsilon\).

Call a family \(\mathscr G\subset\Lcal(H)\) \emph{numerically norming for \(H\)} if
\[
\|x\|^2=\sup_{T\in\mathscr G}|\langle Tx,x\rangle|
\qquad (x\in H).
\]

Set
\[
\eta_0(\varepsilon)=\min\{\varepsilon^2/4,1/4\}\qquad(\varepsilon>0).
\]
Let \(\mathscr U_{\eta_0}(H)\) be the set of all compact operators
\(T\in\Lcal(H)\) such that \(1=\nu(T)<\|T\|\), \(T\) attains its numerical
radius, and the \(\Anu(H)\) condition holds with the same modulus
\(\eta_0\).

\section*{Idea of the Proof}

Proposition 2.8 uses a finite-dimensional identity block and an infinite compact skew-symmetric tail. The skew-symmetric part does not contribute to \(\langle Tx,x\rangle\) over a real Hilbert space; only the identity block contributes. Thus if the identity block is \(F\), then
\[
\langle Tx,x\rangle=\|P_Fx\|^2.
\]
Near numerical-radius attainment forces \(x\) close to \(F\), and normalizing \(P_Fx\) gives an exact numerical-radius attaining point. Since we are free to choose \(F\), taking \(F=\operatorname{span}\{x\}\) gives an operator in the source's class that norms any prescribed vector \(x\).  Since every operator in that class has numerical radius \(1\), no operator can exceed \(\|x\|^2\) on a vector \(x\).

\section*{Theorem}

\begin{theorem}
Let \(H\) be a separable infinite-dimensional real Hilbert space. Then
\(\mathscr U_{\eta_0}(H)\) is numerically norming for \(H\):
\[
\|x\|^2=\sup_{T\in\mathscr U_{\eta_0}(H)}|\langle Tx,x\rangle|
\qquad (x\in H).
\]
More concretely, for every nonzero finite-dimensional subspace \(F\subset H\)
there is \(T_F\in\mathscr U_{\eta_0}(H)\) such that \(T_F|_F=I_F\).  Hence the
same class even norms every finite set of vectors simultaneously after taking
\(F\) to contain that set.
\end{theorem}

\begin{proof}
Fix a nonzero finite-dimensional subspace \(F\subset H\). Since \(H\) is separable and infinite-dimensional, choose an orthonormal decomposition
\[
H=F\oplus \bigoplus_{k=1}^{\infty} \operatorname{span}\{u_k,v_k\}.
\]
Choose real scalars \(\beta_1>1\), \(0<\beta_k<1\) for \(k\ge2\), and \(\beta_k\to0\). Define \(T_F\in\Lcal(H)\) by
\[
T_F f=f\quad(f\in F),\qquad
T_Fu_k=-\beta_k v_k,\qquad
T_Fv_k=\beta_k u_k .
\]
The operator is compact: it is the direct sum of the identity on finite-dimensional \(F\) and \(2\times2\) blocks whose norms \(\beta_k\) tend to zero. Also
\[
\|T_F\|=\max\{1,\sup_k\beta_k\}=\beta_1>1.
\]

Let \(x=f+w\), where \(f\in F\) and \(w\in F^\perp\). The block operator on \(F^\perp\) is skew-symmetric, hence
\[
\langle T_Fx,x\rangle=\|f\|^2+\langle T_Fw,w\rangle=\|f\|^2.
\]
Therefore \(\nu(T_F)\le1\), and equality holds on every unit vector of \(F\). Thus \(\nu(T_F)=1\), and \(T_F\) attains its numerical radius.

We next verify the uniform \(\Anu(H)\) estimate. Let \(0<\varepsilon\le1\), and suppose \(x\in S_H\) satisfies
\[
|\langle T_Fx,x\rangle|>1-\varepsilon^2/4.
\]
Writing \(x=f+w\) as above, we get
\[
\|f\|^2>1-\varepsilon^2/4,\qquad \|w\|<\varepsilon/2.
\]
Set \(x_1=f/\|f\|\in S_F\). Then \(\langle T_Fx_1,x_1\rangle=1\), and
\[
\|x-x_1\|\le \|w\|+\left|1-\|f\|\right|
<\varepsilon/2+\varepsilon^2/4<\varepsilon .
\]
Thus the same modulus \(\eta(\varepsilon)=\varepsilon^2/4\) works for
\(0<\varepsilon\le1\).  For \(\varepsilon>1\), the estimate for
\(\varepsilon=1\) is stronger than required.  Hence \(\eta_0\) is an admissible
common modulus for every \(T_F\).

Since each \(T_F\) is compact, satisfies \(1=\nu(T_F)<\|T_F\|\), attains its
numerical radius, and has common modulus \(\eta_0\), we have
\(T_F\in\mathscr U_{\eta_0}(H)\).  For any \(T\in\mathscr U_{\eta_0}(H)\) and
any \(x\in H\),
\[
|\langle Tx,x\rangle|\le \nu(T)\|x\|^2=\|x\|^2.
\]
Conversely, if \(x\neq0\), take \(F=\operatorname{span}\{x\}\). Then \(T_Fx=x\), so
\[
|\langle T_Fx,x\rangle|=\|x\|^2.
\]
The equality is trivial for \(x=0\). Hence \(\mathscr U_{\eta_0}(H)\) is
numerically norming for \(H\).
\end{proof}

\section*{Verification and Limitations}

The proof is a direct parameterization of the source construction, with the finite identity block allowed to vary. No computational verification is relevant.

This is recorded as a full solution to the extracted source question under the
source-local numerical-radius meaning of ``norming'':
\[
\|x\|^2=\sup_{T\in\mathscr U_{\eta_0}(H)}|\langle Tx,x\rangle|.
\]
This is the interpretation supported by Proposition 2.8, where the relevant
scalar is \(\langle Tx,x\rangle\).  The paper does not define an additional
dual-pairing meaning in which operators would be norming for the Banach space
\(\Lcal(H)\) itself; no claim is made about such a different, external
formulation.

Bounded novelty check: the run indexes had no full packet for arXiv:1910.05726
before this upgrade.  Web searches on June 23, 2026 for the source title, the
phrase ``We do not know how often this happens'', the phrase ``norming for the
whole space'', and related phrases involving \(\Anu\), numerical radius, and
norming did not reveal a later answer. The arXiv abstract page for 1910.05726
was used to confirm source metadata.

\section*{Human Review Recommendation}

Check only the interpretive scope: if ``norming for the whole space'' is read
in the numerical-radius sense used in the source paragraph, the theorem answers
the question.  If a reviewer wants a different operator-space dual-pairing
problem, that should be split off as a separate, newly formalized question.

\begin{thebibliography}{9}

\bibitem{DantasJungRoldan}
S. Dantas, M. Jung, and O. Roldan,
\emph{Norm attaining operators which satisfy a Bollobas type theorem},
arXiv:1910.05726.

\bibitem{BonsallDuncan}
F. F. Bonsall and J. Duncan,
\emph{Numerical Ranges of Operators on Normed Spaces and of Elements of Normed Algebras},
London Mathematical Society Lecture Note Series 2, Cambridge University Press, 1971.

\end{thebibliography}

\end{document}
